% paper/sections/05b_time_integration.tex
% §5 続き：時間積分スキーム（TVD-RK3 + AB2+IPC + CN 法）・安定性制約
%
% ─────────────────────────────────────────────────────────────────────────────
\subsection{時間積分スキーム}
\label{sec:time_int}
% ─────────────────────────────────────────────────────────────────────────────

\noindent
本節では，CLS 移流だけでなく\textbf{NS 方程式系全体}の時間積分フレームワークを扱う．
CLS 移流・再初期化には TVD-RK3（§~\ref{sec:time_int}.1），
NS 運動量方程式には AB2 + IPC + Crank--Nicolson（§~\ref{sec:ab2_ipc}）を用いる．

\subsubsection{3次 TVD Runge--Kutta 法（CLS 移流・再初期化のみ）}

\textbf{CLS 移流ステップ（界面追跡）および CLS 再初期化の仮想時間積分}に，
Shu--Osher 型 TVD-RK3 法~\cite{ShuOsher1988} を用いる：
\begin{align}
  q^{(1)} &= q^n + \Delta t \, \mathcal{L}(q^n) \notag\\
  q^{(2)} &= \frac{3}{4}q^n + \frac{1}{4}q^{(1)} + \frac{1}{4}\Delta t\,\mathcal{L}(q^{(1)}) \notag\\
  q^{n+1} &= \frac{1}{3}q^n + \frac{2}{3}q^{(2)} + \frac{2}{3}\Delta t\,\mathcal{L}(q^{(2)})
  \label{eq:tvd_rk3}
\end{align}
ここで $\mathcal{L}(q)$ は空間演算子であり，
CLS 移流では式~\eqref{eq:dccd_adv_filter} の Dissipative CCD フラックス発散を用いた
$\mathcal{L}(\psi) = -\tilde{F}$（アルゴリズム~\ref{alg:dccd_adv}，負符号は PDE $\partial\psi/\partial t = -\bnabla\cdot(\psi\mathbf{u})$ に由来）を指す．
\textbf{NS 方程式の時間積分には TVD-RK3 は使用しない}（注意~\ref{warn:tvd_rk3_scope}参照）．

\noindent\textbf{TVD（Total Variation Diminishing）の定義：}
数値解の全変動
$TV(q) \equiv \sum_i |q_{i+1} - q_i|$ が時間的に増大しない性質を指す：
$TV(q^{n+1}) \leq TV(q^n)$．
TVD-RK3 はこの性質を時間方向に保証するが，
Dissipative CCD は線形フィルタであり空間 TVD 性を持たない（注意~\ref{warn:adv_clamp}）．
クランプ処置（式~\eqref{eq:psi_clamp}）と組み合わせることで
$\psi \in [0,1]$ を維持する．

\label{warn:tvd_rk3_scope}
TVD-RK3 は以下の2箇所にのみ適用する：
\begin{itemize}
  \item \textbf{CLS 移流方程式}（式~\eqref{eq:cls_advection}）：$\pd{\psi}{t} + \bnabla\cdot(\psi\,\bu) = 0$（保存形）
  \item \textbf{CLS 再初期化方程式}（式~\eqref{eq:cls_reinit_iso}）の仮想時間 $\tau$ 積分（移流部のみ）
\end{itemize}
\textbf{NS 方程式の時間積分}には TVD-RK3 は使用しない（下節 §~\ref{sec:ab2_ipc} 参照）．

\subsubsection{NS 方程式の時間積分：AB2 + IPC 法}
\label{sec:ab2_ipc}

本稿では，プロジェクション法の全体時間精度を $\Ord{\Delta t^2}$ に保つため，
対流項に \textbf{2段 Adams-Bashforth（AB2）法}を採用し，
van Kan（1986）~\cite{vanKan1986} の \textbf{Incremental Pressure Correction（IPC）法}と組み合わせる．
Predictor ステップの離散式は：
\begin{equation}
\begin{split}
  \frac{\rho^{n+1}}{\Delta t}(\bu^* - \bu^n)
  &= -\rho^{n+1}\!\left[\tfrac{3}{2}\mathcal{C}(\bu^n,\bu^n)
    - \tfrac{1}{2}\mathcal{C}(\bu^{n-1},\bu^{n-1})\right] \\
  &\quad + \tfrac{1}{2Re}\bigl[\mathcal{D}(\mu^{n+1},\bu^*)
    + \mathcal{D}(\mu^n,\bu^n)\bigr]
  - \bnabla p^n
  + \bm{f}^{n+1}
\end{split}
  \label{eq:predictor_ab2_ipc}
\end{equation}
ここで $\mathcal{C}(\bu,q) = \bu\cdot\bnabla q$（対流演算子），$\mathcal{D}$ は完全粘性応力演算子
（式~\eqref{eq:viscous_full_x}）．
体積力 $\bm{f}^{n+1}$ は重力と CSF 表面張力の合計で構成する：
\begin{equation*}
  \bm{f}^{n+1} = - \frac{\rho^{n+1}}{Fr^2}\hat{\bm{y}}
  + \frac{1}{\mathrm{We}}\,\kappa^{n+1}\bnabla\psi^{n+1}
\end{equation*}
ここで第1項が重力，第2項が CSF 表面張力 $(\kappa/\mathrm{We})\bnabla\psi$ を表す．
重力は式~\eqref{eq:gravity_convention} の定義 $\bm{g}^* = -g\hat{\bm{z}}$ に対応し，
2次元鉛直問題では $\hat{\bm{z}} = \hat{\bm{y}}$ とする．
圧力場には HFE（§~\ref{sec:field_extension}）を適用し，CCD 勾配が $\Ord{h^6}$ 精度を維持できるよう
界面を跨いで滑らかに延長する（§~\ref{sec:algorithm} 参照）．
対流項の外挿係数 $(3/2,\,-1/2)$ が AB2 であり，
前ステップ圧力 $-\bnabla p^n$ の陽的組み込みが IPC の本質である．

\noindent\textbf{起動ステップ（$n=0$）の注意：}
$n=0$ では $\bu^{n-1}$ が未定義であるため，AB2 外挿係数 $(3/2,\,-1/2)$ の代わりに
前進 Euler 係数 $(1,\,0)$ を用いる
（式~\eqref{eq:predictor_ab2_ipc} の $\mathcal{C}$ 係数を $3/2 \to 1$，$-1/2 \to 0$ と置き換える）．
$\Ord{\Delta t^2}$ の全体時間精度は $n \geq 1$ 以降のステップで達成される．

\paragraph{IPC 法の核心：圧力増分 $\delta p$ の求解}%
\label{sec:ipc_derivation}

IPC 法では，圧力 Poisson 方程式において絶対圧力 $p^{n+1}$ の代わりに
\textbf{圧力増分 $\delta p \equiv p^{n+1} - p^n$} を未知変数として解く
（§~\ref{sec:algorithm} Step 6--7 参照）：
\begin{align*}
  &\text{PPE：}\quad \bnabla\cdot\!\left(\frac{1}{\rho^{n+1}}\bnabla(\delta p)\right)
    = \frac{1}{\Delta t}\bnabla\cdot\bu^*\\
  &\text{補正：}\quad \bu^{n+1} = \bu^* - \frac{\Delta t}{\rho^{n+1}}\bnabla(\delta p),
  \qquad p^{n+1} = p^n + \delta p
\end{align*}
CCD 演算子 $\mathcal{L}^{\rho}_{\mathrm{CCD}}$ は絶対圧力版と\textbf{完全に同一}であり，
求解対象変数が $p^{n+1}$ から $\delta p$ に変わるだけである．
この定式化により，標準 Chorin 法の人工的圧力境界層に起因する
$\Ord{\Delta t}$ スプリッティング誤差が $\Ord{\Delta t^2}$ に低減される
（Guermond et al., \textit{Acta Numerica}, 2006~\cite{Guermond2006}）．

\paragraph{全体的な時間精度}%
\label{sec:time_accuracy_table}

\begin{center}
\renewcommand{\arraystretch}{1.3}
\begin{tabular}{lcc}
\hline
方程式・項 & 時間積分法 & 時間精度 \\
\hline
CLS 移流 & TVD-RK3 & $\Ord{\Delta t^3}$ \\
NS 対流項（Predictor） & AB2（van Kan + 外挿係数 3/2, -1/2） & $\Ord{\Delta t^2}$ \\
NS 粘性項（対角項 $\mathcal{D}_{\mathrm{d}}$） & Crank--Nicolson（半陰的） & $\Ord{\Delta t^2}$ \\
NS 粘性項（クロス微分項 $\mathcal{D}_{\mathrm{c}}$） & 陽的（時刻 $n$） & $\Ord{\Delta t}$（界面近傍）$^\dagger$ \\
Projection スプリッティング & IPC（van Kan 1986） & $\Ord{\Delta t^2}$ \\
\hline
\textbf{全体（均質流域）} & & $\boldsymbol{\Ord{\Delta t^2}}$ \\
\textbf{全体（界面近傍，$\mu_l/\mu_g \gg 1$）} & & $\boldsymbol{\Ord{\Delta t}}$\,$^\dagger$ \\
\hline
\end{tabular}
\end{center}

{\small $^\dagger$ 変粘性界面（$|\bnabla\mu| \neq 0$）ではクロス微分項 $\mathcal{D}_{\mathrm{c}}$ の陽的処理が $\Ord{\Delta t}$ 誤差を導入し，界面近傍での全体時間精度が $\Ord{\Delta t}$ に律速される（注意~\ref{warn:cn_cross_derivative} 参照）．$\mu_l/\mu_g \gg 1$ の場合は AB2 外挿 $[\frac{3}{2}\mathcal{D}_{\mathrm{c}}^n - \frac{1}{2}\mathcal{D}_{\mathrm{c}}^{n-1}]$ の適用を推奨する．}

本スキームの\textbf{均質流域での全体時間精度は $\Ord{\Delta t^2}$（2次精度）}である．
界面近傍（$\mu_l/\mu_g \gg 1$）では上記 $^\dagger$ の注意が適用される．
空間精度（Dissipative CCD 移流：実効 $\Ord{\varepsilon_d^{\mathrm{adv}} h^2}$ フィルタ誤差，
CCD：$\Ord{h^6}$ 微分精度）は，移流フィルタ誤差 $\Ord{\varepsilon_d^{\mathrm{adv}} h^2}$ が律速する
全体空間精度と整合する．CSF モデル誤差 $\Ord{\varepsilon^2}$ が空間精度の律速であり，HFE（§~\ref{sec:field_extension}）によりそれ以外のコンポーネントは $\Ord{h^6}$ を維持する（定量値は §~\ref{sec:bench_droplet} 参照）．

\noindent\textbf{設計上の注意：「CLS は $\Ord{\Delta t^3}$ なのに全体が $\Ord{\Delta t^2}$ な理由」}

CLS 移流に TVD-RK3（$\Ord{\Delta t^3}$）を採用している一方，NS 対流項は AB2（$\Ord{\Delta t^2}$）にとどめている．
この非対称性は意図的な設計選択であり，以下の理由による：

界面と流体運動は\textbf{疎結合（operator-split）}として扱われる．
CLS 移流（界面輸送）は密度場 $\rho^{n+1}$ と曲率 $\kappa^{n+1}$ を決定する「前処理」であり，
その時間精度が $\Ord{\Delta t^2}$ 以上であれば，全体の流体運動精度は NS 方程式側の $\Ord{\Delta t^2}$ で律速する：
\[
  \text{全体時間精度} = \min\!\bigl(\underbrace{\Ord{\Delta t^3}}_\text{CLS},\;
    \underbrace{\Ord{\Delta t^2}}_\text{NS (AB2+IPC)}\bigr) = \Ord{\Delta t^2}
\]
\textbf{CLS の $\Ord{\Delta t^3}$ 高精度化の役割：}
時間精度を $\Ord{\Delta t^3}$ に向上させることには貢献しないが，
界面プロファイルの変形が小さく保たれることで
（i）曲率 $\kappa$ の数値精度が向上し寄生流れを抑制，
（ii）再初期化の必要頻度が減少し質量保存誤差が低減する，
という実用上の優位が得られる．

\noindent\textbf{背景：標準 Chorin 法のスプリッティング誤差（なぜ AB2 だけでは不十分か）}

対流項を AB2（2次精度）に変更するだけでは全体精度は $\Ord{\Delta t^2}$ に\textbf{ならない}．
標準的な Predictor--Corrector 分割（Chorin, 1968~\cite{Chorin1968}）は，
圧力境界条件の不整合（人工的圧力境界層）に起因する $\Ord{\Delta t}$ スプリッティング誤差を独立に内包するため，
対流項の高次化だけで律速が解消されない．
IPC 法による予測ステップへの $p^n$ の組み込みがスプリッティング誤差を除去する唯一の手段である．

\begin{tcolorbox}[warnbox, title={AB2 スタートアップ問題：初回ステップの特別処理}]
\label{warn:ab2_startup}

AB2 法は2時刻分の対流項履歴
$\bigl(\mathcal{C}^n,\,\mathcal{C}^{n-1}\bigr)$
を必要とするため，$t=0$ の最初のステップ（$n=0\to1$）では AB2 を適用できない．

\begin{center}
\renewcommand{\arraystretch}{1.3}
\begin{tabular}{lll}
\hline
ステップ & 対流項の扱い & 時間精度 \\
\hline
$n=0\to1$（初回）& 前進 Euler：$-\mathcal{C}(\bu^0,\bu^0)$ & $\Ord{\Delta t}$（スタートアップ） \\
$n=1\to2$ 以降 & AB2：$\tfrac{3}{2}\mathcal{C}^n - \tfrac{1}{2}\mathcal{C}^{n-1}$ & $\Ord{\Delta t^2}$ \\
\hline
\end{tabular}
\end{center}

初回ステップの前進 Euler による局所打ち切り誤差は $\Ord{\Delta t^2}$（1次精度スキームを1ステップのみ適用した標準値）であり，
以後の AB2 ステップへの伝播が\textbf{安定}であれば，全体積分精度 $\Ord{\Delta t^2}$ に影響しない．

\noindent\textbf{AB2 の零安定性による正当化：}
AB2 は線形多段スキームであり，特性多項式 $\rho(z) = z^2 - z$ の根は $z=0,\,z=1$ で，
いずれも $|z| \leq 1$ かつ $z=1$ は単根（Dahlquist の根条件を満足）．
これにより AB2 は\textbf{零安定（zero-stable）}であり，
$n=0$ での局所誤差 $e_1 = \Ord{\Delta t^2}$ は各後続ステップで高々定数倍（$\leq C$，$C$ は有界な定数）にしか増幅されない
（\cite{DahlquistBjorck1974} §5.3）．
$N_t = T/\Delta t$ ステップ積分で，初回誤差の寄与は $\leq C\,\Ord{\Delta t^2}$（$N_t$ 倍されない）；
残り $N_t - 1$ ステップの AB2 局所誤差合算は $\Ord{(N_t-1)\Delta t^3} = \Ord{T\Delta t^2}$ であり，
全体精度は $\Ord{\Delta t^2}$ に保たれる．

初回完了後，対流項 $\mathcal{C}^0 = \mathcal{C}(\bu^0,\bu^0)$ をバッファに保存し，
以降の AB2 ステップで利用する（追加メモリ：速度場1つ分，$2N_xN_y$ 実数値）．
\end{tcolorbox}

\subsubsection{Crank--Nicolson 法（粘性項）}

粘性項（2階空間微分を含む）は，純陽的処理では厳しい時間刻み制限（$\Delta t < O(h^2)$）を課す．
これを緩和するため，Crank--Nicolson（CN）法による\textbf{半陰的}処理を採用する（行列構造と Thomas 法の詳細は \ref{box:cn_tridiagonal} 参照）：
\begin{equation}
  \frac{u^{n+1} - u^n}{\Delta t} = \frac{1}{2}\bigl[\mathcal{D}(u^n) + \mathcal{D}(u^{n+1})\bigr]
  \label{eq:crank_nicolson}
\end{equation}
ここで $\mathcal{D}$ は粘性拡散演算子．CN 法は時間方向に2次精度（$\Ord{\Delta t^2}$）かつ
\textbf{無条件安定}（スペクトル半径の制約なし）であり，適度な $\Delta t$ での安定計算が可能となる．

\begin{tcolorbox}[warnbox, title={CN 粘性項のクロス微分：$u$ と $v$ の陰的結合と実装上の分離}]
\label{warn:cn_cross_derivative}

変粘性流体の粘性応力 $\bnabla\cdot[\mu(\bnabla\bu + \bnabla^T\bu)]$ の $x$ 成分は：
\begin{align}
  \bigl[\bnabla\cdot(\mu(\bnabla\bu+\bnabla^T\bu))\bigr]_x
  &= \underbrace{2\frac{\partial}{\partial x}\!\left(\mu\frac{\partial u}{\partial x}\right)
     + \frac{\partial}{\partial y}\!\left(\mu\frac{\partial u}{\partial y}\right)}_{\text{対角項 $\mathcal{D}_{\mathrm{d}}(u)$}}
   + \underbrace{\frac{\partial}{\partial y}\!\left(\mu\frac{\partial v}{\partial x}\right)}_{\text{クロス微分項 $\mathcal{D}_{\mathrm{c}}(v)$}}
  \label{eq:viscous_full_x}
\end{align}
クロス微分項 $\mathcal{D}_{\mathrm{c}}(v)$ を CN で $t^{n+1}$ に置くと $u^*$-$v^*$ の連立系（$2N^2$ 未知数）が生じ，ブロック Thomas 法の優位が失われる．
\textbf{標準対処法：クロス項のみ時刻 $n$ で陽的評価し，対角項だけ CN で半陰的処理}する：
\begin{equation*}
  \frac{u^* - u^n}{\Delta t}
  = \frac{1}{2Re}\bigl[\mathcal{D}_{\mathrm{d}}(u^n) + \mathcal{D}_{\mathrm{d}}(u^*)\bigr]
  + \frac{1}{Re}\,\mathcal{D}_{\mathrm{c}}(v^n)
  + \bm{f}_x^{n+1}
\end{equation*}
これにより $u^*$, $v^*$ を独立にブロック Thomas 法で求解できる．

クロス微分項の1次陽的処理は $\Ord{\Delta t}$ 誤差を導入し，気液界面が存在する限り（$|\bnabla\mu| \neq 0$）\textbf{界面近傍での全体時間精度は必ず $\Ord{\Delta t}$ に律速される}．$\mu_l/\mu_g \gg 10$ の水/空気系では特に顕著であり，AB2 外挿 $[\frac{3}{2}\mathcal{D}_{\mathrm{c}}^n - \frac{1}{2}\mathcal{D}_{\mathrm{c}}^{n-1}]$ の適用を推奨する（式~\eqref{eq:predictor_ab2_ipc} の $\mathcal{C}$ 係数と同様の外挿）．
\end{tcolorbox}

\begin{tcolorbox}[warnbox, title={変粘性界面でのクロス微分陽的処理：粘性 CFL 条件の実質的復活}]
\label{warn:cross_cfl}

気液界面近傍では $\mu$ が数格子幅で $\mu_l/\mu_g \approx 10^2$ 変化する．クロス微分 $\mathcal{D}_{\mathrm{c}}(v)$ を陽的処理すると，界面上でのフォン・ノイマン安定性解析から局所 CFL 条件が復活する：
\begin{equation}
  \Delta t \le C_{\mathrm{cross}} \cdot \frac{h^2}{\Delta\mu / \rho}
  \label{eq:cross_cfl}
\end{equation}
ここで $C_{\mathrm{cross}} = O(1)$，$\Delta\mu = \mu_l - \mu_g$．$\mu_l/\mu_g = 100$ の系では等粘性系の陽的粘性 CFL より約 $1/100$ 倍厳しくなる可能性がある．

\begin{center}
\renewcommand{\arraystretch}{1.3}
\begin{tabular}{lll}
\hline
処理方法 & 安定条件 & 適用場面 \\
\hline
対角粘性項 $\mathcal{D}_{\mathrm{d}}$（CN 法） & 無条件安定 & 常に適用 \\
クロス微分項 $\mathcal{D}_{\mathrm{c}}$（陽的） & 式~\eqref{eq:cross_cfl} の制約あり & $\mu_l/\mu_g \lesssim 10$ では無害 \\
\hline
\end{tabular}
\end{center}

\textbf{シミュレーションが爆発した場合：}対流 CFL・毛管波制約（式~\eqref{eq:dt_adv},\eqref{eq:dt_sigma}）を確認後，なお爆発する場合はクロス微分項 $\mathcal{D}_{\mathrm{c}}$ が原因の可能性がある．$\Delta t$ を $1/10$ に減らして安定化するか確認することを推奨する．
\end{tcolorbox}

\begin{tcolorbox}[algbox, title={CN 対角粘性系の行列構造と Thomas 法による求解}]
\label{box:cn_tridiagonal}
対角粘性項 $\mathcal{D}_{\mathrm{d}}(u) = 2\partial_x(\mu\partial_x u) + \partial_y(\mu\partial_y u)$
の CN 半陰的処理で，$x$ 方向項 $(1/2Re)\cdot 2\partial_x(\mu\partial_x u^*)$
をLHSに移項すると（$(1/2Re)\cdot 2 = 1/Re$），行 $j$ の連立系は：
\[
  \left[\frac{\rho^{n+1}_{ij}}{\Delta t}
    + \frac{\mu_{i+1/2,j}+\mu_{i-1/2,j}}{Re\,h^2}\right]u^*_{ij}
  - \frac{\mu_{i+1/2,j}}{Re\,h^2}\,u^*_{i+1,j}
  - \frac{\mu_{i-1/2,j}}{Re\,h^2}\,u^*_{i-1,j}
  = \mathrm{RHS}_{ij}
\]
の\textbf{三重対角形}となる（下・上対角係数は負；$D_{d,x}$ の因子 $2$ と CN の $1/2$ が相殺して $1/Re$）：

\medskip
\textbf{$x$ 方向スウィープ（各行 $j$ 独立）：}
$[\mathbf{L}_x^j]\{u^*_{\cdot,j}\} = \{\mathrm{RHS}^j\}$ を Thomas 法で $\mathcal{O}(N_x)$ 求解．
\begin{itemize}
  \item 主対角：$\rho^{n+1}/\Delta t + (\mu_{i+1/2}+\mu_{i-1/2})/(Re\,h^2)$
  \item 上・下対角：$-\mu_{i+1/2}/(Re\,h^2)$，$-\mu_{i-1/2}/(Re\,h^2)$（ともに負符号）
\end{itemize}

\textbf{$y$ 方向スウィープ（各列 $i$ 独立）：}
$y$ 方向の $D_{d,y}(u) = \partial_y(\mu\partial_y u)$（因子 $1$）に対し：
\begin{itemize}
  \item 上・下対角：$-\mu_{j+1/2}/(2\,Re\,h^2)$（$y$ 方向は $D_{d,y}$ の因子 $1$ と $1/2Re$ でそのまま $1/(2Re)$）
  \item $x$ 方向と異なる係数に注意．
\end{itemize}
Thomas 法で $\mathcal{O}(N_y)$ 求解．

\textbf{境界条件の組み込み：}
壁面（ノースリップ，$u^*=0$）では対応行を単位行に置き換える．

\textbf{$x$-$y$ 分割と全体精度：}
本稿の実装では $y$ 方向 CN を時刻 $n$ の陽的評価で補い（注意~\ref{warn:cn_cross_derivative}），
$u^*$-$v^*$ の連立系を $\mathcal{O}(N)$ Thomas 法で完全分離する
（§~\ref{sec:time_accuracy_table} 参照）．
\end{tcolorbox}
